\chapter{Conclusións}
\label{chap:conclusions}

\lettrine{D}{erradeiro} capítulo da memoria, onde se presentará a
situación final do traballo, as leccións aprendidas, a relación coas
competencias da titulación en xeral e a mención en particular,
posibles liñas futuras,\dots


\section{TFG en sí}
Teño que dala sensación de que traballei dabondo
\section{Resultados}

\begin{itemize}
				\item \nota{o cruce non funciona en absoluto, polo menos en imagenet. o cal non era trivial en absoluto. e ademáis pola teoría ate parecería o contrario.}
				\item \nota{a Rede que comparte ofrece mellor separación semantica ca normal no dataset RFMiD}
				\item \nota{crate nunca usouse para cousas medicas -> isto é a parte de INTERPRETABILIDADE do meu TFG. se non con DINO valdría}
				\item \nota{en canto a rendemento puro parece que a mellor variante é a orixinal. non obstante para estar seguros é preciso máis experimentación.}
				\item \nota{en canto a mapas de atención parece que en Imagenet a variante segundo orde aínda con peor rendemento ofrecía mellor separación semantica (falta probas cuantitativas).}
				\item \nota{en canto a palabras do dicionario non teño conclusión ningunha. :/}
				\item \nota{o segundo orde parece que precisa un lr máis baixo. ten sentido.}
				\item \nota{en xeral pouca conclusion-conclusion teño, penso que máis que conclusions o que teño son moitas más preguntas}
\end{itemize}


\section{Traballo futuro}

\begin{itemize}
				\item \nota{repetir o experimento de imagenet con máis épocas? e todas as variantes?}
				\item \nota{repetir o experimento de rfmid con máis épocas? e todas as variantes? e clasificación multiclase? e incorporando MAO?}
				\item \nota{extender o principio de compartir $U$ e $D$ á matriz $W$? e ós batchnorm? co fin de ter unha rede totalmente repetida? explorar ata cando pódese compartir?}
\end{itemize}
